\documentclass[11pt]{article}
\usepackage{amsmath,amssymb,amsthm,algorithm,algorithmic}
\usepackage{booktabs}
\usepackage[margin=1in]{geometry}

\newtheorem{theorem}{Theorem}
\newtheorem{lemma}[theorem]{Lemma}
\newtheorem{corollary}[theorem]{Corollary}
\newtheorem{proposition}[theorem]{Proposition}
\newtheorem{definition}[theorem]{Definition}

\title{Problem 2: Even Fibonacci Numbers}
\author{Project Euler Solutions}
\date{}

\begin{document}
\maketitle

\section{Problem Statement}

By considering the terms in the Fibonacci sequence whose values do not exceed four million, find the sum of the even-valued terms. That is, compute
\[
S = \sum_{\substack{F_n \le 4\,000\,000 \\ 2 \mid F_n}} F_n
\]
where $F_1 = 1$, $F_2 = 1$, and $F_n = F_{n-1} + F_{n-2}$ for $n \ge 3$.

\section{Mathematical Development}

\subsection{Preliminaries}

\begin{definition}[Fibonacci Sequence]
The Fibonacci sequence $(F_n)_{n \ge 1}$ is defined by $F_1 = 1$, $F_2 = 1$, and $F_n = F_{n-1} + F_{n-2}$ for $n \ge 3$.
\end{definition}

\begin{lemma}[Fibonacci Addition Identity]
\label{lem:add}
For all integers $m \ge 2$ and $n \ge 1$,
\[
F_{m+n} = F_m F_{n+1} + F_{m-1} F_n.
\]
\end{lemma}

\begin{proof}
Fix $m \ge 2$ and proceed by induction on $n$.

\emph{Base case.} $n = 1$: $F_{m+1} = F_m F_2 + F_{m-1} F_1 = F_m + F_{m-1}$, which is the defining recurrence.

\emph{Inductive step.} Assume the identity holds for $n$ and $n-1$ (with $n \ge 2$). Then
\begin{align*}
F_{m+n+1} &= F_{m+n} + F_{m+n-1} \\
           &= (F_m F_{n+1} + F_{m-1} F_n) + (F_m F_n + F_{m-1} F_{n-1}) \\
           &= F_m(F_{n+1} + F_n) + F_{m-1}(F_n + F_{n-1}) \\
           &= F_m F_{n+2} + F_{m-1} F_{n+1}. \qedhere
\end{align*}
\end{proof}

\subsection{Parity Periodicity}

\begin{theorem}[Pisano Period Modulo 2]
\label{thm:parity}
$F_n$ is even if and only if $3 \mid n$.
\end{theorem}

\begin{proof}
The recurrence $F_n \equiv F_{n-1} + F_{n-2} \pmod{2}$ depends only on the pair $(F_{n-1} \bmod 2,\, F_{n-2} \bmod 2)$. Since there are at most $2^2 = 4$ distinct pairs, the sequence of pairs is eventually periodic. Computing:
\[
(F_n \bmod 2)_{n \ge 1} = (1,\, 1,\, 0,\, 1,\, 1,\, 0,\, \ldots)
\]
The pair $(F_4 \bmod 2,\, F_5 \bmod 2) = (1,1) = (F_1 \bmod 2,\, F_2 \bmod 2)$, so the parity sequence is periodic with period~3. Since $F_1 \bmod 2 = 1 \ne 0 = F_3 \bmod 2$, the minimal period is exactly~3. Hence $2 \mid F_n$ iff $3 \mid n$.
\end{proof}

\subsection{The Even Fibonacci Recurrence}

\begin{definition}
Let $E_k = F_{3k}$ for $k \ge 1$. By Theorem~\ref{thm:parity}, $(E_k)$ is the subsequence of all even Fibonacci numbers.
\end{definition}

\begin{theorem}[Even Fibonacci Recurrence]
\label{thm:recurrence}
For all $k \ge 3$, $E_k = 4E_{k-1} + E_{k-2}$, with $E_1 = 2$ and $E_2 = 8$.
\end{theorem}

\begin{proof}
We show $F_{3k} = 4F_{3k-3} + F_{3k-6}$.

Applying Lemma~\ref{lem:add} with $m=3$, $n=3k-3$:
\begin{equation}
F_{3k} = F_3 F_{3k-2} + F_2 F_{3k-3} = 2F_{3k-2} + F_{3k-3}. \tag{1}
\end{equation}

Since $F_{3k-2} = F_{3k-3} + F_{3k-4}$:
\begin{equation}
F_{3k} = 3F_{3k-3} + 2F_{3k-4}. \tag{2}
\end{equation}

Applying Lemma~\ref{lem:add} with $m=3$, $n=3k-6$ gives $F_{3k-3} = 2F_{3k-5} + F_{3k-6}$, so $F_{3k-5} = (F_{3k-3} - F_{3k-6})/2$. Since $F_{3k-4} = F_{3k-5} + F_{3k-6}$:
\begin{equation}
F_{3k-4} = \frac{F_{3k-3} + F_{3k-6}}{2}. \tag{3}
\end{equation}

Substituting (3) into (2):
\[
F_{3k} = 3F_{3k-3} + (F_{3k-3} + F_{3k-6}) = 4F_{3k-3} + F_{3k-6}.
\]

Initial conditions: $E_1 = F_3 = 2$, $E_2 = F_6 = 8$. Verification: $E_3 = F_9 = 34 = 4 \cdot 8 + 2$.
\end{proof}

\subsection{Growth Rate and Closed Form}

\begin{theorem}[Closed Form]
\label{thm:closed}
\[
E_k = \frac{\alpha^k - \beta^k}{\sqrt{5}}
\]
where $\alpha = 2 + \sqrt{5}$ and $\beta = 2 - \sqrt{5}$ are roots of $x^2 - 4x - 1 = 0$.
\end{theorem}

\begin{proof}
The characteristic polynomial of $E_k = 4E_{k-1} + E_{k-2}$ is $x^2 - 4x - 1 = 0$ with roots $\alpha = 2+\sqrt{5}$ and $\beta = 2 - \sqrt{5}$. The general solution is $E_k = c_1 \alpha^k + c_2 \beta^k$. From initial conditions and using $\alpha\beta = -1$, $\alpha + \beta = 4$, $\alpha - \beta = 2\sqrt{5}$:

From $E_1 = c_1\alpha + c_2\beta = 2$ and $E_2 = c_1\alpha^2 + c_2\beta^2 = 8$. Since $\alpha^2 = 4\alpha + 1$ and $\beta^2 = 4\beta + 1$, we get $E_2 = 4E_1 + (c_1 + c_2) = 8 + (c_1 + c_2)$, so $c_1 + c_2 = 0$. Then $E_1 = c_1(\alpha - \beta) = 2c_1\sqrt{5} = 2$, giving $c_1 = 1/\sqrt{5}$, $c_2 = -1/\sqrt{5}$.
\end{proof}

\begin{corollary}[Exponential Growth]
$E_k = \Theta(\alpha^k)$ where $\alpha = 2 + \sqrt{5} \approx 4.236$. Equivalently, $\alpha = \phi^3$ where $\phi = (1+\sqrt{5})/2$.
\end{corollary}

\begin{proof}
$|\beta| = |2-\sqrt{5}| < 1$, so $|\beta^k| \to 0$. Thus $E_k \sim \alpha^k/\sqrt{5}$. That $\alpha = \phi^3$ follows from $\phi^3 = \phi \cdot \phi^2 = \phi(\phi+1) = \phi^2 + \phi = 2\phi + 1 = 2 + \sqrt{5}$.
\end{proof}

\subsection{Summation Formula}

\begin{theorem}[Partial Sum Identity]
\label{thm:sum}
$\displaystyle T_k = \sum_{j=1}^{k} E_j = \frac{E_{k+1} + E_k - 2}{4}$.
\end{theorem}

\begin{proof}
By induction on $k$.

\emph{Base case.} $T_1 = 2 = (8+2-2)/4$. \checkmark

\emph{Inductive step.} Assume $T_{k-1} = (E_k + E_{k-1} - 2)/4$. Then
\[
T_k = T_{k-1} + E_k = \frac{E_k + E_{k-1} - 2 + 4E_k}{4} = \frac{5E_k + E_{k-1} - 2}{4}.
\]
We need $5E_k + E_{k-1} = E_{k+1} + E_k$, i.e., $E_{k+1} = 4E_k + E_{k-1}$, which is the recurrence of Theorem~\ref{thm:recurrence}.
\end{proof}

\subsection{Numerical Evaluation}

\begin{proposition}
$S = 4\,613\,732$.
\end{proposition}

\begin{proof}
The even Fibonacci terms not exceeding $4\,000\,000$:
\begin{center}
\begin{tabular}{crr}
\toprule
$k$ & $E_k$ & $\sum_{j=1}^{k} E_j$ \\
\midrule
1  & 2 & 2 \\
2  & 8 & 10 \\
3  & 34 & 44 \\
4  & 144 & 188 \\
5  & 610 & 798 \\
6  & 2\,584 & 3\,382 \\
7  & 10\,946 & 14\,328 \\
8  & 46\,368 & 60\,696 \\
9  & 196\,418 & 257\,114 \\
10 & 832\,040 & 1\,089\,154 \\
11 & 3\,524\,578 & 4\,613\,732 \\
\bottomrule
\end{tabular}
\end{center}
Since $E_{12} = 4 \cdot 3\,524\,578 + 832\,040 = 14\,930\,352 > 4\,000\,000$, the sum includes $k = 1, \ldots, 11$.

Verification via Theorem~\ref{thm:sum}: $T_{11} = (14\,930\,352 + 3\,524\,578 - 2)/4 = 18\,454\,928/4 = 4\,613\,732$.
\end{proof}

\section{Algorithm}

\begin{algorithm}[H]
\caption{\textsc{SumEvenFibonacci}$(L)$}
\begin{algorithmic}[1]
\STATE $a \gets 2$, $b \gets 8$, $\mathrm{total} \gets 0$
\WHILE{$a \le L$}
    \STATE $\mathrm{total} \gets \mathrm{total} + a$
    \STATE $(a, b) \gets (b,\; 4b + a)$
\ENDWHILE
\RETURN total
\end{algorithmic}
\end{algorithm}

\begin{theorem}[Correctness]
\textsc{SumEvenFibonacci}$(L)$ returns $\sum\{E_k : E_k \le L\}$ for all $L \ge 1$.
\end{theorem}

\begin{proof}
Loop invariants at the start of iteration $i$: (1)~$a = E_i$, $b = E_{i+1}$ (by Theorem~\ref{thm:recurrence}); (2)~$\mathrm{total} = \sum_{j=1}^{i-1} E_j$. The loop terminates when $E_i > L$ (guaranteed since $E_k \to \infty$), at which point $\mathrm{total} = \sum_{j=1}^{i-1} E_j = \sum\{E_k : E_k \le L\}$.
\end{proof}

\section{Complexity Analysis}

\begin{theorem}
\textsc{SumEvenFibonacci}$(L)$ runs in $\Theta(\log L)$ time and $O(1)$ space.
\end{theorem}

\begin{proof}
Each iteration performs $O(1)$ work. The loop executes $K = \max\{k : E_k \le L\}$ iterations. Since $E_k = \Theta(\alpha^k)$ with $\alpha > 1$, we have $K = \Theta(\log_\alpha L) = \Theta(\log L)$. The algorithm stores 3 integer variables, using $O(1)$ space.
\end{proof}

\section{Answer}

\[
\boxed{4613732}
\]

\end{document}
